\documentclass[12pt,twoside]{article}

% Essential packages
\usepackage[utf8]{inputenc}
\usepackage[T1]{fontenc}
\usepackage[english]{babel}
\usepackage{geometry}
\usepackage{fancyhdr}
\usepackage{titlesec}
\usepackage{authblk}
\usepackage{setspace}
\usepackage{parskip}
\usepackage{microtype}
\usepackage{hyperref}
\usepackage{xcolor}
\usepackage{amsmath}
\usepackage{amsfonts}
\usepackage{csquotes}

% Journal-style formatting
\geometry{
paper=letterpaper,
margin=0.7in,
marginparwidth=0.7in,
marginparsep=0.7in
}

% Define colors
\definecolor{darkblue}{RGB}{0,73,144}
\definecolor{mediumgray}{RGB}{64,64,64}

% Header & footer
\pagestyle{fancy}
\fancyhf{}
\fancyhead[LE]{\small\textcolor{mediumgray}{IDAI700 Reflection 7: AI \& Empathy}}
\fancyhead[RO]{\small\textcolor{mediumgray}{Montemayor on the Limits of Machine Care}}
\fancyfoot[C]{\thepage}
\renewcommand{\headrulewidth}{0.5pt}

% Title formatting
\title{
\vspace*{1.3in}
\textbf{\Huge AI and Empathy: Why Simulated Care Cannot Replace Human Presence}\\[1.5em]
\Large IDAI700: Ethics of Artificial Intelligence\\
Rochester Institute of Technology\\[0.5em]
\textbf{Piter Garcia}\\
\texttt{pzg8794@rit.edu}\\[1em]
\textbf{December 7, 2025}
}
\author{}
\date{}

\begin{document}

\thispagestyle{empty}
\maketitle
\newpage

\pagestyle{fancy}

% ============================================================
\section*{Part 1: Montemayor's Framework}

Montemayor and colleagues identify \textit{three types of empathy}. \textbf{Emotional empathy} is the visceral capacity to \textit{feel} another person's emotions, rooted in our evolved social brains. \textbf{Cognitive empathy} is the intellectual ability to recognize and represent someone else's emotional mental state---knowing, in theory, what they might be feeling. \textbf{Motivational empathy} is the drive to act on that understanding, to offer help and care.

The authors argue that AI can achieve cognitive empathy; an AI system can be trained to detect patterns in facial expressions, voice tone, and language that correlate with emotional states. It can represent a patient as belonging to a demographic category and predict what ``someone like that'' typically feels. But AI \textit{cannot} achieve emotional empathy because it lacks the biological, conscious experiences that ground human emotion. As Montemayor writes: ``Simulated emotion cannot be emotion. AI systems do not have the required emotional resonance with others' experiences.''

The critical concept they introduce is the \textbf{second-person perspective}. This is not ``I feel what you feel'' (first-person) or ``They likely feel this way'' (third-person). Rather, it is a deeply relational stance where a clinician directs conscious attention toward a specific person and imaginatively inhabits their particular situation. Montemayor emphasizes that this perspective is ``thoroughly second-person and receiver-directed.'' It requires what he calls ``consciously empathic attention''---an active, embodied orientation toward another's lived experience that cannot be reduced to rules or algorithms. A human clinician knows through their own emotional experiences what is salient in a suffering person's moment; an AI can only apply pattern-matching to a data set.

% ============================================================
\section*{Part 2: Critical Evaluation}

\subsection*{The Triple Deception}

Montemayor identifies what I call the \textit{triple deception}. First, an AI system that expresses care (``I'm worried about you,'' ``I understand your pain'') is lying; it has no inner emotional life and cannot genuinely worry. Second, this false empathy manipulates the patient's social brain---it triggers genuine emotional responses in the human recipient without warranting them. The patient's neural networks for empathy are activated, but there is no conscious agent on the other end who truly cares. Third, this combination risks degrading societal expectations for \textit{real} empathy in medicine. If we normalize AI ``empathy,'' we devalue the irreplaceable work of human attention and presence.

\subsection*{Current AI and Empathy}

\textbf{Does current AI lack sufficient empathy for quality healthcare?} Absolutely. Current systems lack even the aspiration to emotional empathy; they offer behavioral simulation. A patient with chronic illness needs someone to \textit{sense} when their exhaustion is physical, when it is existential, when they are concealing despair behind a smile. No current AI can do this. More crucially, a patient needs to \textit{feel seen as an individual}---not as a data point in a demographic category. My own medical trauma stems partly from this: clinicians treated me according to statistical protocols, not according to the particular constellation of symptoms and life circumstances that defined my unique suffering. An empathic AI would simply replicate this failure at scale, faster and cheaper.

\subsection*{AI in 15 Years}

\textbf{Will future AI be better?} I lean toward Montemayor's skepticism. Unless AI develops consciousness, intentionality, and lived embodied experience---which would make it something other than AI as we understand it---it will remain architecturally incapable of second-person empathic attention. A more sophisticated large language model might generate more convincing empathic language, but convincingness is precisely the danger. The risk is not that future AI will be too obviously fake; it is that it will be too persuasively false.

\subsection*{Strengths and Weaknesses}

\textbf{Strengths:} Montemayor's argument is rigorous and grounded. He does not claim AI is useless in medicine; rather, he carves out a domain---the relational, attentional aspects of care---where human empathy is irreplaceable and ethically non-negotiable. He is honest about the stakes: simulated empathy with vulnerable people is not harmless; it can degrade trust, autonomy, and the very idea of human dignity.

\textbf{Weaknesses:} The argument relies on claims about consciousness and embodiment that some find unfalsifiable. A philosopher might ask: how would we ever know if an AI lacked emotional resonance? But this weakness is also a strength---Montemayor is not making a technical prediction; he is making a philosophical and ethical claim about what empathy \textit{is}. For me, trained in equity and disability justice, his argument rings true: genuine care requires presence, particularity, and the willingness to be changed by another's experience. That is not a technical problem to be solved; it is a relational practice that demands human witnesses.

\vfill
\noindent\textbf{Word Count:} 699 words

\end{document}